\chapter{Пакетный аудит гипотез: четыре леммы и новые прунеры}

Десять гипотез были повторно проверены в замороженных конвенциях
Протокола~0.1: train-якоря \(\{\alpha^{-1},m_e,m_\mu\}\), единичные
радиусы, спектральное меню
\(\eta_D(\mathbb{RP}^3)=\pm\tfrac14\) и окружной функционал
\begin{equation}
 \mathcal I_\rho(\beta)=
 \log\frac{\cosh(2\pi\rho)-\cos(2\pi\beta)}
 {\cosh(2\pi\rho)-1}.
\end{equation}
Аудит не добавил закрытого физического предсказания, но дал несколько
точных уточнений и сузил пространство повторных поисков.

\section{Воспроизведённые отрицательные результаты}

Для APS-кандидата необходимо различать две допустимые ветви. При
\(\eta=-1/4\)
\begin{equation}
 -\frac{\eta}{2}+|\zeta_R(-1)|
 =\frac18+\frac1{12}=\frac5{24},
\end{equation}
а при \(\eta=+1/4\) та же комбинация равна \(-1/24\). Ни одна ветвь не
даёт \(1/3\), а значение \(\eta=-1/2\) отсутствует в замороженном меню.
Независимо от этой арифметики вектороподобный блок
\begin{equation}
 (m+iE)(m-iE)=m^2+E^2
\end{equation}
сокращает фазу по каждому уровню, поэтому \(\eta\)-часть не создаёт
аддитивного вещественного массового сдвига.

Конформный и диффеоморфный фактор действует только в метрическом блоке.
Прямая сумма с майорановским пространством оставляет его ранг равным
\(24\), так что маршрут \(24-1=23\) вновь закрывается отрицательно.

Матрица
\begin{equation}
 \operatorname{diag}(-1,-1,-1,+1,+1)
\end{equation}
имеет определитель \(-1\) и не принадлежит \(SU(5)\). Исправленная
\(\mathbb Z_2/\mathbb Z_4\)-проекция с
\begin{equation}
 P_5=\operatorname{diag}(1,1,1,-1,-1),\qquad
 h=\operatorname{diag}(-1,-1,-1,-i,-i)
\end{equation}
точно оставляет \(U+2D+H\) и направление
\((17/6,1/6,2)\), но не выводит величину порога и не исправляет знак
обычного промежуточного ренормгруппового пробега.

Предложенный интегранд
\(\operatorname{Tr}(A\wedge dA+\tfrac23A^3)\) является трёхформой, а не
пятиформой. Корректный шаблон имеет вид
\begin{equation}
 \omega_5(A)=\operatorname{Tr}
 \left(AF^2-\frac12A^3F+\frac1{10}A^5\right),
 \qquad d\omega_5=\operatorname{Tr}F^3,
\end{equation}
с точностью до конвенции и точных форм. Это исправляет размерность, но не
создаёт отсутствующий операторный мост от
\(\eta_D(\mathbb{RP}^3)\) к границе
\(\mathbb{RP}^3\times S^1\).

\section{Сохранённые точные леммы}

Обычная шестиканальная гауссова сумма даёт
\begin{equation}
 \operatorname{Var}(\bar x)
 =\frac1{6\zeta(4,1/2)}=\frac1{\pi^4},
\end{equation}
но нормированный след даёт \(6/\pi^4\), а полный физический носитель и
нормировка источника не выведены.

Минимальный \(U(1)_{B-L}\)-контент с одним \(N^c\) на поколение проходит
все шесть непрерывных аномальных сумм, даёт корневую фазу \(i\) и
калибровочно-инвариантную вершину поля заряда \(-2\). Однородный конденсат
закрыт трилеммой, поэтому сохраняется только неоднородная текстура.
Смешанный след
\begin{equation}
 \sum_{\mathrm{Weyl}}Y(B-L)=\frac83
\end{equation}
задаёт будущий нормировочно-чувствительный второй гейт, но ещё не является
вторым пройденным сектором.

Для приведённого седлового анзаца точно сохраняются
\begin{equation}
 k_n=\frac{2\pi n+\pi}{L},\qquad
 \lambda v^2>\frac{\pi^2}{L^2},
\end{equation}
а при \(L=\pi\) --- порог \(\lambda v^2>1\) и вырожденная пара
\(n=0,-1\). Само превышение порога и выбор ориентации действием не
зафиксированы.

\section{Новые точные прунеры}

Окружной функционал чётен по голономной ветви:
\begin{equation}
 \mathcal I_\rho(\beta)=\mathcal I_\rho(-\beta)
 =\mathcal I_\rho(1-\beta).
\end{equation}
Следовательно, он тождественно не расщепляет сопряжённые ориентации.
Численное сравнение при \(\rho=1\),
\begin{align}
 \mathcal I_1(1/2)&=0.0074697796100663\ldots,\\
 2\mathcal I_1(1/4)&=0.0074837289794912\ldots,
\end{align}
является отдельным неравенством, а не доказательством вырождения; их
отношение равно \(0.9981360403\ldots\).

Попытка заменить член \(-\alpha/3\) двухпетлевой формой
\begin{equation}
 \left(\frac{\alpha}{\pi}\right)^2X=\frac\alpha3
\end{equation}
требует
\begin{equation}
 X=\frac{\pi^2}{3\alpha}=450.8303668617\ldots.
\end{equation}
Строгий дефект этой гипотезы состоит не в субъективной
``неестественности'' числа, а в циркулярности: коэффициент получен из
train-якоря \(\alpha\), который он должен был независимо объяснить.

Кандидат проекционного веса
\begin{equation}
 \frac{12\pi}{7}=5.3855874062\ldots
\end{equation}
не равен требуемому
\(5.4027533072\ldots\); относительный дефицит равен
\(3.17725\times10^{-3}\). Без независимого оператора это приближение
отклоняется Протоколом~0.1.

\section{Незарегистрированное утверждение}

Фраза ``осевой разностный сдвиг тождественно равен нулю'' не сопровождается
в пакете определением соответствующего оператора. Существующий аудит осей
доказывает одну остаточную проективную орбиту, а относительный двухсекторный
аудит оставляет свободный вещественный угол смешивания. Поэтому новый
нулевой операторный прунер не регистрируется до предъявления явной формулы.

\begin{proposition}[Итог пакетного аудита]
Сохраняются четыре точных или условно точных конструктивных результата:
шестиканальная дисперсия, \(B-L\)-аномальная свобода и корневая фаза,
неоднородный порог с сопряжённой парой и исправленная
\(\mathbb Z_2/\mathbb Z_4\)-проекция представлений. Нового закрытого
физического предсказания не возникает:
\begin{equation}
 N_{\mathrm{closed\ physical}}=0.
\end{equation}
Следующий допустимый шаг требует заранее фиксированного родительского
действия, а не новой вставки, восстановленной из целевого числа.
\end{proposition}

Машинный протокол сохранён в
\texttt{s2t\_hypothesis\_batch\_pruner\_results.json}.